\subsection{Experiment \rnum{1}: Model training and Reconstruction}

\begin{table}[!htbp]
    \centering
    \setlength{\tabcolsep}{10pt}
    \begin{tabular}{l cccc}
        \toprule
        \rowcolor{gray!20}
        \textbf{Metric} & \textbf{AE} & \textbf{$\beta$-VAE} & \textbf{CAE} & \textbf{$\beta$-CVAE} \\
        \midrule
        Best Train loss & 3.075e-4 & 4.592e-3 & \textbf{1.428e-4} & 4.712e-3 \\
        \rowcolor{gray!10} Best Validation loss & 2.941e-4 & 4.592e-3 & \textbf{4.669e-5} & 4.592e-3 \\
        Epochs before best loss & \textbf{3} & 33 & 10 & 33 \\
        \rowcolor{gray!10} Model Sizes (F16) & 5.12 GB & 5.13 GB & \textbf{91.4 KB} & 126.53 MB \\
        F32 Inference speed (ms) & \textbf{1.658} & 654.14 & 267.31 & 979.52 \\
        \rowcolor{gray!10} F16 Inference speed (ms) & \textbf{1.875} & 6.31 & 5.671 & \textbf{1.857} \\
        \bottomrule
    \end{tabular}
    \caption{Training comparison between the model.}
    \label{tab:modelresinfo}
    \smallskip
    \small{*The reconstruction error is the sum of all reconstruction errors across all batches.}
\end{table}
%\begin{figure}[!h]
%    \centering
%    \includegraphics[scale=0.35]{figures/time.png}
%    \caption{Comparison between Maximum and mediant amount of time spent per epoch for training and validation combined. The maximum only happens on first epoch, while the median times stay within a 3\% margin for the other epochs.}
%    \label{fig:traintimes}
%\end{figure}

Figure \ref{fig:traintimes} compares training times, training, and validation steps combined across all models. The max time only happens on the first epoch, while every subsequent epoch lasts within a 5\% range of the median time. We notice a considerable speedup between the first and consecutive epochs for the non-variational autoencoders. However, the variational models achieve a lower initial training time. 

\begin{figure}[htbp]
\centering
\begin{tikzpicture}
\begin{axis}[
    ybar,
    bar width=15pt,
    width=\textwidth,
    height=0.6\textwidth,
    legend style={at={(0.5,-0.15)},
                  anchor=north,legend columns=-1},
    ylabel={Time (minutes)},
    symbolic x coords={AE,$\beta$-VAE,CAE,$\beta$-CVAE},
    xtick=data,
    nodes near coords,
    nodes near coords align={vertical},
    ymin=0,ymax=60,
    xlabel={Models},
    enlarge x limits=0.2,
]
\addplot[fill=lightgreen!80] coordinates {(AE,41.6) ($\beta$-VAE,18.95) (CAE,45.73) ($\beta$-CVAE,18.95)}; 
\addplot[fill=red!80] coordinates {(AE,3.60) ($\beta$-VAE,15.42) (CAE,10.45) ($\beta$-CVAE,15.42)}; 
\legend{Max Time, Mean Time}
% Add text above each group of bars
\node [above] at (axis cs:AE,60) {41m 36s};
\node [above] at (axis cs:$\beta$-VAE,60) {18m 57s};
\node [above] at (axis cs:CAE,60) {45m 44s};
\node [above] at (axis cs:$\beta$-CVAE,60) {18m 57s};
\end{axis}
\end{tikzpicture}
\caption{Comparison of Max and Mean Time (Train + Validation) per epoch}
\label{fig:traintimes}
\end{figure}
\clearpage
\subsubsection{Train and Validation Losses}
\begin{figure}[!htbp]
  \centering
  \begin{subfigure}{\textwidth}
    \centering
    \includegraphics[width=0.7\textwidth]{figures/losses/ae.png}
    \caption{AE: We observe a rapid loss, followed by no big change in either train or validation loss, indicating that the model has learned, but is unable to }
  \end{subfigure}
  
  \vspace{0.5cm}
  
  \begin{subfigure}{\textwidth}
    \centering
    \includegraphics[width=0.7\textwidth]{figures/losses/vae.png}
    \caption{$\beta$-VAE: The loss rapidly declines before flattening out. Around epoch 15, a sudden drop is observed, indicating a shift in learning rate, which allows the loss to further decline}
  \end{subfigure}
  
  \vspace{0.5cm}
  
  \begin{subfigure}{\textwidth}
    \centering
    \includegraphics[width=0.7\textwidth]{figures/losses/cae.png}
    \caption{CAE: Initially, a more gradual decline is observed. After epoch 4, the validation loss starts to shift, before stabilizing around epoch 12}
  \end{subfigure}
  
  \vspace{0.5cm}
  
  \begin{subfigure}{\textwidth}
    \centering
    \includegraphics[width=0.7\textwidth]{figures/losses/cvae.png}
    \caption{$\beta$-CVAE. A linear decrease is observed, with the validation loss expected to be lower than the training loss.}
  \end{subfigure}
  
  \caption{Train and Validation Loss over Time. Train loss is blue, and validation loss is orange}
  \label{fig:losses}
\end{figure}

\clearpage

\subsubsection{Half-precision Reconstruction Heatmaps}
% Page 1: Original Heatmaps
\begin{figure}[!p]
\centering
\begin{subfigure}{\textwidth}
    \centering
    \includegraphics[width=0.65\textwidth]{figures/anomalies/before/20190415_031735.png}
    \caption{2019-04-15 03:17:35}
\end{subfigure}
\vspace{1em}
\begin{subfigure}{\textwidth}
    \centering
    \includegraphics[width=0.65\textwidth]{figures/anomalies/before/20190415_031750.png}
    \caption{2019-04-15 03:17:50}
\end{subfigure}
\vspace{1em}
\begin{subfigure}{\textwidth}
    \centering
    \includegraphics[width=0.65\textwidth]{figures/anomalies/before/20190415_031755.png}
    \caption{2019-04-15 03:17:55}
\end{subfigure}
\caption{Original DAS heatmaps at 3 different timestamps}
\label{fig:original_heatmaps}
\end{figure}
\clearpage
% Page 2: AE Reconstructions
\begin{figure}[p]
\centering
\begin{subfigure}{\textwidth}
    \centering
    \includegraphics[width=0.9\textwidth]{figures/anomalies/ae/20190415_031735.png}
    \caption{2019-04-15 03:17:35}
\end{subfigure}
\vspace{1em}
\begin{subfigure}{\textwidth}
    \centering
    \includegraphics[width=0.9\textwidth]{figures/anomalies/ae/20190415_031750.png}
    \caption{2019-04-15 03:17:50}
\end{subfigure}
\vspace{1em}
\begin{subfigure}{\textwidth}
    \centering
    \includegraphics[width=0.9\textwidth]{figures/anomalies/ae/20190415_031755.png}
    \caption{2019-04-15 03:17:55}
\end{subfigure}
\caption{AE reconstructions}
\label{fig:ae_reconstructions}
\end{figure}
\clearpage
% Page 3: β-VAE Reconstructions
\begin{figure}[p]
\centering
\begin{subfigure}{\textwidth}
    \centering
    \includegraphics[width=0.9\textwidth]{figures/anomalies/vae/20190415_031735.png}
    \caption{2019-04-15 03:17:35}
\end{subfigure}
\vspace{1em}
\begin{subfigure}{\textwidth}
    \centering
    \includegraphics[width=0.9\textwidth]{figures/anomalies/vae/20190415_031750.png}
    \caption{2019-04-15 03:17:50}
\end{subfigure}
\vspace{1em}
\begin{subfigure}{\textwidth}
    \centering
    \includegraphics[width=0.9\textwidth]{figures/anomalies/vae/20190415_031755.png}
    \caption{2019-04-15 03:17:55}
\end{subfigure}
\caption{$\beta$-VAE reconstructions}
\label{fig:beta_vae_reconstructions}
\end{figure}
\clearpage
% Page 4: CAE Reconstructions
\begin{figure}[p]
\centering
\begin{subfigure}{\textwidth}
    \centering
    \includegraphics[width=0.9\textwidth]{figures/anomalies/cae/20190415_031735.png}
    \caption{2019-04-15 03:17:35}
\end{subfigure}
\vspace{1em}
\begin{subfigure}{\textwidth}
    \centering
    \includegraphics[width=0.9\textwidth]{figures/anomalies/cae/20190415_031750.png}
    \caption{2019-04-15 03:17:50}
\end{subfigure}
\vspace{1em}
\begin{subfigure}{\textwidth}
    \centering
    \includegraphics[width=0.9\textwidth]{figures/anomalies/cae/20190415_031755.png}
    \caption{2019-04-15 03:17:55}
\end{subfigure}
\caption{CAE reconstructions}
\label{fig:cae_reconstructions}
\end{figure}
\clearpage
% Page 5: β-CVAE Reconstructions
\begin{figure}[p]
\centering
\begin{subfigure}{\textwidth}
    \centering
    \includegraphics[width=0.9\textwidth]{figures/anomalies/cvae/20190415_031735.png}
    \caption{2019-04-15 03:17:35}
\end{subfigure}
\vspace{1em}
\begin{subfigure}{\textwidth}
    \centering
    \includegraphics[width=0.9\textwidth]{figures/anomalies/cvae/20190415_031750.png}
    \caption{2019-04-15 03:17:50}
\end{subfigure}
\vspace{1em}
\begin{subfigure}{\textwidth}
    \centering
    \includegraphics[width=0.9\textwidth]{figures/anomalies/cvae/20190415_031755.png}
    \caption{2019-04-15 03:17:55}
\end{subfigure}
\caption{$\beta$-CVAE reconstructions}
\label{fig:beta_cvae_reconstructions}
\end{figure}